\chapter{\eh{globe-showing-americas} Introducción a las Redes de Datos}

%% ── 1.1 ¿Qué es una red? ─────────────────────────────────────────────────
\section{\eh{thinking-face} ¿Qué es una red de datos?}

\begin{concepto}
Una \textbf{red de datos} es un conjunto de dispositivos (computadoras, teléfonos,
impresoras, servidores) conectados entre sí para poder \textbf{compartir información
y recursos}.

\emoji{light-bulb} \textbf{Analogía:} Imagina que tienes amigos en diferentes casas.
Una red es como un sistema de tuberías que conecta todas las casas para que
puedan mandarse agua (datos) entre sí. Cada casa tiene su dirección (IP), y
las tuberías (cables o Wi-Fi) transportan el agua hasta la puerta correcta.
\end{concepto}

\section{\eh{books} Un poco de historia \emoji{hourglass-done}}

\begin{itemize}
    \item \textbf{1969 --- ARPANET:} La primera red de computadoras del mundo,
          financiada por el Departamento de Defensa de EE.UU. Solo 4 nodos:
          UCLA, Stanford, UCSB y University of Utah.
          \emoji{exploding-head} El primer mensaje enviado fue \textit{``lo''}
          porque el sistema crasheó antes de terminar de enviar \textit{``login''}!
    \item \textbf{1974 --- TCP/IP:} Vint Cerf y Bob Kahn publican el paper que
          define los protocolos de Internet.
    \item \textbf{1983 --- Nacimiento de Internet:} ARPANET adopta TCP/IP
          oficialmente. Se le llama el ``Big Bang'' de Internet.
    \item \textbf{1991 --- World Wide Web:} Tim Berners-Lee inventa HTTP y HTML
          en el CERN. Internet ya no es solo para militares e investigadores.
    \item \textbf{2003 --- Wi-Fi se masifica:} 802.11g llega a laptops de consumo.
    \item \textbf{2015 --- HTTP/2:} Google impulsa la segunda versión de HTTP.
    \item \textbf{2021 --- QUIC (RFC 9000):} Google revoluciona el transporte web.
          Hoy $\approx$25\% del tráfico de Internet usa QUIC.
    \item \textbf{2022 --- HTTP/3:} HTTP sobre QUIC estandarizado (RFC 9114).
\end{itemize}

%% ── 1.3 Modelo OSI ───────────────────────────────────────────────────────
\section{\eh{birthday-cake} El Modelo OSI --- Las 7 Capas}

\begin{concepto}
El \textbf{modelo OSI} (Open Systems Interconnection) es un modelo de referencia
creado por ISO en 1984. Divide la comunicación en red en \textbf{7 capas},
cada una con una responsabilidad específica y bien definida.

\emoji{light-bulb} \textbf{Analogía:} Es como enviar una carta por DHL. Tú
escribes el mensaje (capa 7), lo metes en un sobre (capa 6), le pones dirección
y número de seguimiento (capas 3--4), el empaquetador lo mete en una caja (capa 2),
el camión lo transporta por carretera (capa 1). Cada quien hace su parte
sin saber exactamente qué hacen los demás.
\end{concepto}

\begin{center}
\begin{tikzpicture}[
    layer/.style={rectangle, draw, minimum width=9cm, minimum height=0.85cm,
                  font=\bfseries\small, align=center, rounded corners=2pt},
]
\node[layer, fill=red!25]    (L7) at (0,5.1) {Capa 7 --- \emoji{desktop-computer} Aplicación (Application)};
\node[layer, fill=orange!30] (L6) at (0,4.25){Capa 6 --- \emoji{artist-palette} Presentación (Presentation)};
\node[layer, fill=yellow!35] (L5) at (0,3.4) {Capa 5 --- \emoji{handshake} Sesión (Session)};
\node[layer, fill=green!25]  (L4) at (0,2.55){Capa 4 --- \emoji{truck} Transporte (Transport)};
\node[layer, fill=cyan!25]   (L3) at (0,1.7) {Capa 3 --- \emoji{world-map} Red (Network)};
\node[layer, fill=blue!20]   (L2) at (0,0.85){Capa 2 --- \emoji{link} Enlace de Datos (Data Link)};
\node[layer, fill=purple!20] (L1) at (0,0.0) {Capa 1 --- \emoji{electric-plug} Física (Physical)};
\draw[-{Stealth}, thick] (L7.south) -- (L6.north);
\draw[-{Stealth}, thick] (L6.south) -- (L5.north);
\draw[-{Stealth}, thick] (L5.south) -- (L4.north);
\draw[-{Stealth}, thick] (L4.south) -- (L3.north);
\draw[-{Stealth}, thick] (L3.south) -- (L2.north);
\draw[-{Stealth}, thick] (L2.south) -- (L1.north);
\end{tikzpicture}
\end{center}

\begin{table}[h!]
\centering
\renewcommand{\arraystretch}{1.3}
\begin{tabular}{clll}
\toprule
\textbf{Capa} & \textbf{Nombre} & \textbf{Qué hace} & \textbf{Ejemplos} \\
\midrule
7 & Aplicación   & Interfaz con el usuario y apps       & HTTP, DNS, SMTP, FTP \\
6 & Presentación & Codifica, comprime y cifra           & TLS, Base64, JPEG \\
5 & Sesión       & Gestiona sesiones de comunicación    & NetBIOS, RPC \\
4 & Transporte   & Entrega extremo a extremo, puertos   & TCP, UDP, QUIC \\
3 & Red          & Enrutamiento entre redes distintas   & IP, ICMP, OSPF, BGP \\
2 & Enlace       & Entrega en la misma red local (LAN)  & Ethernet, MAC, ARP \\
1 & Física       & Bits sobre el medio físico           & UTP, Fibra, Radio \\
\bottomrule
\end{tabular}
\caption{Resumen del modelo OSI}
\end{table}

\begin{cuidado}
\textbf{\emoji{brain} Mnemónica (capas 7 a 1):}

\textbf{A}ll \textbf{P}eople \textbf{S}eem \textbf{T}o \textbf{N}eed \textbf{D}ata \textbf{P}rocessing

(Application, Presentation, Session, Transport, Network, Data Link, Physical)
\end{cuidado}

%% ── 1.4 Modelo TCP/IP ─────────────────────────────────────────────────────
\section{\eh{laptop} El Modelo TCP/IP (Modelo de Internet)}

\begin{concepto}
El \textbf{modelo TCP/IP} es el que usa Internet en la realidad. Fue creado
por DARPA en los 70s y tiene \textbf{4 capas}. No es un modelo teórico sino
el protocolo real que funciona hoy en cada dispositivo conectado a Internet.
\end{concepto}

\begin{table}[h!]
\centering
\renewcommand{\arraystretch}{1.4}
\begin{tabular}{lll}
\toprule
\textbf{Capa TCP/IP} & \textbf{Equivale a OSI} & \textbf{Protocolos} \\
\midrule
Aplicación    & Capas 7 + 6 + 5 & HTTP, DNS, SMTP, SSH, FTP \\
Transporte    & Capa 4          & TCP, UDP, QUIC \\
Internet      & Capa 3          & IPv4, IPv6, ICMP \\
Acceso a red  & Capas 2 + 1     & Ethernet, Wi-Fi, ARP \\
\bottomrule
\end{tabular}
\caption{Modelo TCP/IP vs OSI}
\end{table}

%% ── 1.5 PDUs ──────────────────────────────────────────────────────────────
\section{\eh{package} PDUs --- ¿Cómo se llaman los datos en cada capa?}

Cada capa llama de distinta forma a la unidad de datos que maneja
(\textbf{PDU}: Protocol Data Unit):

\begin{table}[h!]
\centering
\renewcommand{\arraystretch}{1.3}
\begin{tabular}{clll}
\toprule
\textbf{Capa} & \textbf{PDU} & \textbf{Contiene} & \textbf{Analogía} \\
\midrule
7--5 & Datos (Data)            & El mensaje de la app    & El texto de la carta \\
4    & Segmento / Datagrama    & Datos + puertos         & Carta con remitente y destinatario \\
3    & Paquete (Packet)        & Segmento + IPs          & Caja con dirección postal \\
2    & Trama (Frame)           & Paquete + MACs          & Caja subida al camión \\
1    & Bits                    & 0s y 1s en el cable     & Las ruedas del camión girando \\
\bottomrule
\end{tabular}
\caption{PDUs por capa OSI}
\end{table}

\begin{ejemplo}
\textbf{Qué pasa cuando abres \texttt{google.com}:}
\begin{enumerate}
    \item Tu browser (capa 7) genera un \textbf{HTTP GET /}
    \item TCP (capa 4) añade puertos (src: 54321, dst: 443) \textrightarrow{} \textbf{segmento}
    \item IP (capa 3) añade tu IP y la de Google \textrightarrow{} \textbf{paquete}
    \item Ethernet (capa 2) añade MACs \textrightarrow{} \textbf{trama}
    \item Tu NIC (capa 1) convierte en \textbf{señal eléctrica} o \textbf{radio}
    \item En Google, el proceso se invierte: reconstruyen el HTTP GET desde los bits
\end{enumerate}
\end{ejemplo}

%% ── 1.6 RFC e IETF ────────────────────────────────────────────────────────
\section{\eh{scroll} RFCs --- Las reglas de Internet}

\begin{concepto}
Un \textbf{RFC} (Request for Comments) es un documento técnico que define cómo
funciona un protocolo de Internet. Son publicados por el \textbf{IETF}
(Internet Engineering Task Force) y son de acceso público y gratuito en
\texttt{https://www.rfc-editor.org}.

\emoji{light-bulb} \textbf{Analogía:} Los RFCs son como las reglas del fútbol
publicadas por la FIFA. Sin reglas comunes, un router Cisco no podría
hablar con uno Juniper, y tu celular no podría conectarse a Internet.
\end{concepto}

\begin{table}[h!]
\centering
\renewcommand{\arraystretch}{1.3}
\begin{tabular}{lll}
\toprule
\textbf{RFC} & \textbf{Protocolo} & \textbf{Año} \\
\midrule
RFC 768  & UDP                           & 1980 \\
RFC 791  & IPv4                          & 1981 \\
RFC 793  & TCP                           & 1981 \\
RFC 1918 & Direcciones IP privadas       & 1996 \\
RFC 2328 & OSPF versión 2                & 1998 \\
RFC 4251 & SSH Architecture              & 2006 \\
RFC 8200 & IPv6                          & 2017 \\
RFC 8446 & TLS 1.3                       & 2018 \\
RFC 8484 & DNS over HTTPS (DoH)          & 2018 \\
RFC 9000 & \textbf{QUIC}                 & \textbf{2021} \\
RFC 9114 & \textbf{HTTP/3}               & \textbf{2022} \\
RFC 9250 & DNS over QUIC (DoQ)           & 2022 \\
\bottomrule
\end{tabular}
\caption{RFCs clave}
\end{table}

\section{\eh{office-building} Organismos de Estandarización}

\begin{table}[h!]
\centering
\renewcommand{\arraystretch}{1.3}
\begin{tabular}{lll}
\toprule
\textbf{Organismo} & \textbf{Nombre completo} & \textbf{Qué estandariza} \\
\midrule
IETF  & Internet Engineering Task Force          & Protocolos de Internet (RFCs) \\
IEEE  & Inst. of Electrical and Electronics Eng. & Ethernet (802.3), Wi-Fi (802.11) \\
ISO   & Intl. Organization for Standardization  & Modelo OSI, ISO 27000 (seguridad) \\
IANA  & Internet Assigned Numbers Authority      & IPs, puertos, nombres de dominio \\
ANSI  & American National Standards Institute    & Estándares industriales EE.UU. \\
TIA   & Telecommunications Industry Association  & Cableado estructurado (con EIA) \\
LACNIC & Latin America and Caribbean NIC        & Asignación de IPs en Latinoamérica \\
\bottomrule
\end{tabular}
\caption{Organismos de estandarización}
\end{table}
